\documentclass[10pt,a4paper]{amsart}
\usepackage{amsmath,amssymb,amsthm,graphicx}
\usepackage{verbatim}
\newenvironment{reference}{\comment}{\endcomment}
%\newenvironment{reference}{}{}
\newenvironment{slogan}{\comment}{\endcomment}
\newenvironment{history}{\comment}{\endcomment}

% For commutative diagrams we use Xy-pic
\usepackage[all]{xy}

% We use 2cell for 2-commutative diagrams.
\xyoption{2cell}
% Theorem environments.
%
\theoremstyle{plain}
\newtheorem{theorem}[subsection]{Theorem}
\newtheorem{proposition}[subsection]{Proposition}
\newtheorem{lemma}[subsection]{Lemma}

\theoremstyle{definition}
\newtheorem{definition}[subsection]{Definition}
\newtheorem{example}[subsection]{Example}
\newtheorem{exercise}[subsection]{Exercise}
\newtheorem{situation}[subsection]{Situation}

\theoremstyle{remark}
\newtheorem{remark}[subsection]{Remark}
\newtheorem{remarks}[subsection]{Remarks}

\numberwithin{equation}{subsection}

% Macros
%
\def\lim{\mathop{\mathrm{lim}}\nolimits}
\def\colim{\mathop{\mathrm{colim}}\nolimits}
\def\Spec{\mathop{\mathrm{Spec}}}
\def\Hom{\mathop{\mathrm{Hom}}\nolimits}
\def\Ext{\mathop{\mathrm{Ext}}\nolimits}
\def\SheafHom{\mathop{\mathcal{H}\!\mathit{om}}\nolimits}
\def\SheafExt{\mathop{\mathcal{E}\!\mathit{xt}}\nolimits}
\def\Sch{\mathit{Sch}}
\def\Mor{\mathop{\mathrm{Mor}}\nolimits}
\def\Ob{\mathop{\mathrm{Ob}}\nolimits}
\def\Sh{\mathop{\mathit{Sh}}\nolimits}
\def\NL{\mathop{N\!L}\nolimits}
\def\CH{\mathop{\mathrm{CH}}\nolimits}
\def\proetale{{pro\text{-}\acute{e}tale}}
\def\etale{{\acute{e}tale}}
\def\QCoh{\mathit{QCoh}}
\def\Ker{\mathop{\mathrm{Ker}}}
\def\Im{\mathop{\mathrm{Im}}}
\def\Coker{\mathop{\mathrm{Coker}}}
\def\Coim{\mathop{\mathrm{Coim}}}

% Boxtimes
%
\DeclareMathSymbol{\boxtimes}{\mathbin}{AMSa}{"02}

%
% Macros for moduli stacks/spaces
%
\def\QCohstack{\mathcal{QC}\!\mathit{oh}}
\def\Cohstack{\mathcal{C}\!\mathit{oh}}
\def\Spacesstack{\mathcal{S}\!\mathit{paces}}
\def\Quotfunctor{\mathrm{Quot}}
\def\Hilbfunctor{\mathrm{Hilb}}
\def\Curvesstack{\mathcal{C}\!\mathit{urves}}
\def\Polarizedstack{\mathcal{P}\!\mathit{olarized}}
\def\Complexesstack{\mathcal{C}\!\mathit{omplexes}}
% \Pic is the operator that assigns to X its picard group, usage \Pic(X)
% \Picardstack_{X/B} denotes the Picard stack of X over B
% \Picardfunctor_{X/B} denotes the Picard functor of X over B
\def\Pic{\mathop{\mathrm{Pic}}\nolimits}
\def\Picardstack{\mathcal{P}\!\mathit{ic}}
\def\Picardfunctor{\mathrm{Pic}}
\def\Deformationcategory{\mathcal{D}\!\mathit{ef}}
\title{Functoriality of the sheaf of differentials for a commutative diagram of ringed topoi}
\author{The Stacks Project Authors}
\date{stacks-project:sites-modules.tex:local-snapshot-20261012; GNU FDL 1.2 or later}
\begin{document}
\maketitle
\noindent\textit{Editorial note.} All delivered mathematical text is the author's own, copied verbatim from the pinned Stacks Project chapter file sites-modules.tex; no statement or proof step was rewritten, reordered or removed.\par\medskip
\noindent\textit{Editorial note.} Both the statement and the proof contain commutative diagrams typeset with Xy-pic; the retained two-slice preamble of this source supplies the native Xy-pic setup, so they compile as in the source.\par\medskip
\noindent\textit{Editorial note.} The counted span cites no labelled external item, so no navigation stubs were needed.\par\medskip
\noindent\textit{Allowed standard prerequisites (editorial).}\begin{itemize}
\item Sheaves of differentials on ringed topoi and their universal derivation. Stacks Project, Modules on Sites, section on differentials (definition of Omega and d\_\{X/Y\}). Supplies the objects and the universal derivation whose functoriality is being established.
\item Adjointness of f\_* and f\textasciicircum{}* for a morphism of ringed topoi. Stacks Project, Modules on Sites, section on morphisms of ringed topoi. Converts the map into the pushforward form into the map c\_f in the pullback form.
\end{itemize}\medskip
\begin{lemma}
\label{lemma-functoriality-differentials-ringed-topoi}
Let
$X = (\Sh(\mathcal{C}_X), \mathcal{O}_X)$,
$Y = (\Sh(\mathcal{C}_Y), \mathcal{O}_Y)$,
$X' = (\Sh(\mathcal{C}_{X'}), \mathcal{O}_{X'})$, and
$Y' = (\Sh(\mathcal{C}_{Y'}), \mathcal{O}_{Y'})$ be ringed topoi.
Let
$$
\xymatrix{
X' \ar[d] \ar[r]_f & X \ar[d] \\
Y' \ar[r] & Y
}
$$
be a commutative diagram of morphisms of ringed topoi. The map
$f^\sharp : \mathcal{O}_X \to f_*\mathcal{O}_{X'}$ composed with the map
$f_*\text{d}_{X'/Y'} : f_*\mathcal{O}_{X'} \to f_*\Omega_{X'/Y'}$ is a
$Y$-derivation. Hence we obtain a canonical map of $\mathcal{O}_X$-modules
$\Omega_{X/Y} \to f_*\Omega_{X'/Y'}$, and by
adjointness of $f_*$ and $f^*$ a
canonical $\mathcal{O}_{X'}$-module homomorphism
$$
c_f : f^*\Omega_{X/Y} \longrightarrow \Omega_{X'/Y'}.
$$
It is uniquely characterized by the property that
$f^*\text{d}_{X/Y}(t)$ maps to $\text{d}_{X'/Y'}(f^* t)$
for any local section $t$ of $\mathcal{O}_X$.
\end{lemma}

\begin{proof}
This is clear except for the last assertion. Let us explain the meaning of
this. Let $U \in \Ob(\mathcal{C}_X)$ and let $t \in \mathcal{O}_X(U)$.
This is what it means for $t$ to be a local section of $\mathcal{O}_X$.
Now, we may think of $t$ as a map of sheaves of sets
$t : h_U^\# \to \mathcal{O}_X$. Then
$f^{-1}t : f^{-1}h_U^\# \to f^{-1}\mathcal{O}_X$. By $f^*t$ we mean
the composition
$$
\xymatrix{
f^{-1}h_U^\# \ar[rr]^{f^{-1}t} \ar@/^4ex/[rrrr]^{f^*t} & &
f^{-1}\mathcal{O}_X \ar[rr]^{f^\sharp} & &
\mathcal{O}_{X'}
}
$$
Note that $\text{d}_{X/Y}(t) \in \Omega_{X/Y}(U)$. Hence we may think of
$\text{d}_{X/Y}(t)$ as a map $\text{d}_{X/Y}(t) : h_U^\# \to \Omega_{X/Y}$.
Then $f^{-1}\text{d}_{X/Y}(t) : f^{-1}h_U^\# \to f^{-1}\Omega_{X/Y}$.
By $f^*\text{d}_{X/Y}(t)$ we mean the composition
$$
\xymatrix{
f^{-1}h_U^\#
\ar[rr]^{f^{-1}\text{d}_{X/Y}(t)}
\ar@/^4ex/[rrrr]^{f^*\text{d}_{X/Y}(t)} & &
f^{-1}\Omega_{X/Y} \ar[rr]^{1 \otimes \text{id}} & &
f^*\Omega_{X/Y}
}
$$
OK, and now the statement of the lemma means that we have
$$
c_f \circ f^*t = f^*\text{d}_{X/Y}(t)
$$
as maps from $f^{-1}h_U^\#$ to $\Omega_{X'/Y'}$. We omit the verification
that this property holds for $c_f$ as defined in the lemma. (Hint: The first
map $c'_f : \Omega_{X/Y} \to f_*\Omega_{X'/Y'}$ satisfies
$c'_f(\text{d}_{X/Y}(t)) = f_*\text{d}_{X'/Y'}(f^\sharp(t))$ as sections of
$f_*\Omega_{X'/Y'}$ over $U$, and you have to
turn this into the equality above by using adjunction.)
The reason that this uniquely characterizes $c_f$ is that the images
of $f^*\text{d}_{X/Y}(t)$ generate the $\mathcal{O}_{X'}$-module
$f^*\Omega_{X/Y}$ simply because the local sections $\text{d}_{X/Y}(t)$
generate the $\mathcal{O}_X$-module $\Omega_{X/Y}$.
\end{proof}

\end{document}
